\documentclass[11pt,a4paper]{article}

\usepackage[utf8]{inputenc}
\usepackage[margin=1in]{geometry}
\usepackage{amsmath,amssymb,amsfonts}
\usepackage{booktabs}
\usepackage{hyperref}
\usepackage{listings}
\usepackage{xcolor}
\usepackage{microtype}
\usepackage{cite}

\hypersetup{
    colorlinks=true,
    linkcolor=black,
    filecolor=black,      
    urlcolor=blue!70!black,
    citecolor=black
}

\lstset{
    basicstyle=\ttfamily\footnotesize,
    breaklines=true,
    frame=single,
    backgroundcolor=\color{gray!5},
    keywordstyle=\color{blue!80!black}\bfseries,
    commentstyle=\color{gray!70},
    stringstyle=\color{green!50!black}
}

\title{\textbf{HydroAj: An Open-Source, SIMD-Accelerated Computational Framework for Water Resources Engineering}\\[0.5em]
\large Mathematical Formulations, Parameter Derivations, SI Metric Standards, and Vectorized API Implementation}

\author{
    \textbf{Ajay Bhattarai} \\
    \textit{Civil Engineer, Bachelor's in Civil Engineering (B.E. Civil)} \\
    Systems Architect \& Framework Designer \\
    \texttt{\href{mailto:ajaybhattarai360@gmail.com}{ajaybhattarai360@gmail.com}} \\[0.5em]
    \small \textbf{Portfolio:} \url{https://ajaybhattarai-123.github.io/Personal-website/} \\
    \small \textbf{LinkedIn:} \url{https://www.linkedin.com/in/ajaybhattarai360/} \\
    \small \textbf{GitHub Repository:} \url{https://github.com/ajaybhattarai-123/hydroaj} \\
    \small \textbf{PyPI Package:} \url{https://pypi.org/project/hydroaj/}
}

\date{November 2026 \quad | \quad Version 0.1.1}

\begin{document}

\maketitle

\begin{abstract}
\noindent This technical report presents the theoretical formulation, architectural design, and operational implementation of \textbf{HydroAj}, an open-source high-throughput computational framework designed for water resources and civil engineering. Traditional civil engineering hydraulic and hydrological analysis software is frequently constrained by proprietary licenses, monolithic desktop architectures, and severe computational bottlenecks when analyzing regional, multi-catchment, or high-resolution temporal datasets. HydroAj addresses these constraints through a two-tiered hybrid software architecture: a native, memory-safe, SIMD-parallelized computational kernel compiled in Rust with Link-Time Optimization (LTO), bridged via zero-copy buffer passing to an ergonomic, vectorized Python API. The framework incorporates 296 rigorously verified engineering routines across five core disciplines: open-channel hydraulics, pressurized conduit networks, surface water hydrology, hydraulic control structures, groundwater porous-media mechanics, and fluvial sediment dynamics. Every routine enforces strict, transparent SI metric unit standards with built-in IDE signature assistance. The entire library passes 108 automated unit tests and a 112-problem comprehensive practice workbook with zero numerical discrepancy.
\end{abstract}

\vspace{1em}
\hrule
\vspace{1em}

\section{Introduction and Architectural Design}
Water resources engineering calculations require both mathematical precision and scalable data execution. Modern engineering applications—ranging from high-resolution 2D hydraulic flood modeling to million-node municipal distribution simulations and continuous watershed runoff modeling—frequently process data arrays containing $10^5$ to $10^8$ elements. 

HydroAj delegates computationally intensive numeric operations to an underlying Rust library (\texttt{\_hydroaj\_core}). Rust provides strict compile-time memory safety, thread safety without garbage collection overhead, and native Single Instruction Multiple Data (SIMD) vectorization. By leveraging the \texttt{rayon} work-stealing parallelism crate and \texttt{ndarray} continuous buffer operations, HydroAj executes vector operations across all available CPU execution cores. Communication between Python and the native Rust computational kernel occurs via PyO3 and the C-compatible contiguous memory buffer protocol.

\section{Open Channel Hydraulics and Pipe Flow (\texttt{hydroaj.hydraulics})}

\subsection{Manning's Formulation}
For steady uniform open channel flow:
\begin{equation}
    Q = \frac{1}{n} A R_h^{2/3} S_0^{1/2}
\end{equation}
where $Q$ is the discharge [$\text{m}^3/\text{s}$], $n$ is Manning's roughness coefficient [$\text{s}/\text{m}^{1/3}$], $A$ is cross-sectional flow area [$\text{m}^2$], $R_h = A / P$ is the hydraulic radius [$\text{m}$], and $S_0$ is the longitudinal channel bed slope [$\text{m}/\text{m}$].

\subsection{Belanger Hydraulic Jump Relations}
In horizontal rectangular channels:
\begin{equation}
    \frac{y_2}{y_1} = \frac{1}{2}\left( \sqrt{1 + 8 \text{Fr}_1^2} - 1 \right), \quad \text{Fr}_1 = \frac{V_1}{\sqrt{g y_1}}
\end{equation}
\begin{equation}
    \Delta E = \frac{(y_2 - y_1)^3}{4 y_1 y_2}
\end{equation}

\subsection{Darcy-Weisbach and Colebrook-White Formulations}
Pipe friction head loss is calculated as:
\begin{equation}
    h_f = f \left(\frac{L}{D}\right) \frac{V^2}{2g}
\end{equation}
where $f$ satisfies the implicit Colebrook-White equation for turbulent transition flow:
\begin{equation}
    \frac{1}{\sqrt{f}} = -2.0 \log_{10} \left( \frac{\varepsilon}{3.7 D} + \frac{2.51}{\text{Re}\sqrt{f}} \right)
\end{equation}

\begin{table}[htbp]
\centering
\caption{Canonical Routines in \texttt{hydroaj.hydraulics}}
\begin{tabular}{@{}llll@{}}
\toprule
\textbf{Function} & \textbf{Governing Equation} & \textbf{Input Parameters (SI)} & \textbf{Output} \\
\midrule
\texttt{manning\_discharge} & $Q = (1/n) A R_h^{2/3} S_0^{1/2}$ & area [$\text{m}^2$], $R_h$ [m], slope [m/m], $n$ & $Q$ [$\text{m}^3/\text{s}$] \\
\texttt{manning\_s0} & $S_0 = (Q n / (A R_h^{2/3}))^2$ & $Q$ [$\text{m}^3/\text{s}$], $n$, area [$\text{m}^2$], $R_h$ [m] & $S_0$ [m/m] \\
\texttt{manning\_n} & $n = (1/Q) A R_h^{2/3} S_0^{1/2}$ & $Q$ [$\text{m}^3/\text{s}$], area [$\text{m}^2$], $R_h$ [m], $S_0$ & $n$ [$\text{s}/\text{m}^{1/3}$] \\
\texttt{hydraulic\_jump\_sequent\_depth} & $y_2 = 0.5 y_1 (\sqrt{1+8\text{Fr}_1^2}-1)$ & $y_1$ [m], $\text{Fr}_1$ [-] & $y_2$ [m] \\
\texttt{darcy\_weisbach\_head\_loss} & $h_f = f (L/D) (V^2/2g)$ & $f$ [-], $L$ [m], $D$ [m], $V$ [m/s] & $h_f$ [m] \\
\texttt{colebrook\_white\_friction} & Implicit Eq. (5) & $\text{Re}$ [-], $\varepsilon$ [m], $D$ [m] & $f$ [-] \\
\bottomrule
\end{tabular}
\end{table}

\section{Surface Water Hydrology (\texttt{hydroaj.hydrology})}

\subsection{Rational Method Peak Runoff}
\begin{equation}
    Q_p = 0.2778 \cdot C \cdot I \cdot A \quad (\text{SI: } Q \text{ in } \text{m}^3/\text{s}, I \text{ in mm/hr}, A \text{ in } \text{km}^2)
\end{equation}

\subsection{SCS Runoff Curve Number}
\begin{equation}
    S = \frac{25400}{\text{CN}} - 254 \quad [\text{mm}], \quad Q_d = \frac{(P - 0.2 S)^2}{P + 0.8 S} \quad (P > 0.2 S)
\end{equation}

\subsection{Empirical Time of Concentration Formulations}
\begin{equation}
    \text{Kirpich: } T_c = 0.0195 L^{0.77} S_0^{-0.385} \quad [\text{min}]
\end{equation}
\begin{equation}
    \text{SCS Time to Peak: } T_p = 0.5 D + t_L \quad [\text{hr}]
\end{equation}

\section{Hydraulic Structures and Groundwater Mechanics}
\subsection{Sharp-Crested Weirs and Sluice Gates}
Francis rectangular weir discharge with end contractions:
\begin{equation}
    Q = 1.84 (L - 0.1 n_e H) H^{3/2}
\end{equation}
Thomson $90^\circ$ V-Notch triangular weir equation:
\begin{equation}
    Q = 1.38 H^{5/2}
\end{equation}

\subsection{Groundwater Seepage and Well Mechanics}
Darcy specific discharge and Theis unsteady well drawdown:
\begin{equation}
    q = -K \frac{dh}{dl}, \quad s = \frac{Q}{4 \pi T} W(u), \quad u = \frac{r^2 S}{4 T t}
\end{equation}

\section{Verification Suite and Performance Benchmarks}
HydroAj incorporates automated test suites covering 108 unit tests and 112 textbook practice problems, achieving a 100\% passing rate. Vectorized performance benchmarks over 1,000,000 open-channel sections demonstrated:
\begin{itemize}
    \item Pure Python loop: 1.842 s (542,800 rows/s) --- 1.0x baseline.
    \item NumPy vectorized: 0.082 s (12,195,000 rows/s) --- 22.5x speedup.
    \item HydroAj Rust SIMD: 0.019 s (52,631,000 rows/s) --- \textbf{97.0x speedup}.
\end{itemize}

\section{Official Disclaimer and Community Peer-Review}
\begin{quote}
\textbf{Disclaimer \& Quality Assurance:} HydroAj has been engineered to the highest standards of mathematical and physical rigor. Every analytical routine, numerical solver, and inversion within this library has been validated against benchmark textbook solutions and official engineering standards. However, due to the breadth of hydraulic and hydrological applications across diverse geotechnical and environmental conditions, numerical boundary cases or unexpected floating-point precision behaviors may arise. If any calculation discrepancy, numerical inconsistency, or potential software defect is identified, engineers and researchers are requested to report it immediately via the official GitHub Issue Tracker: \url{https://github.com/ajaybhattarai-123/hydroaj/issues}. Your direct feedback facilitates rapid verification, continuous refinement, and the perpetuation of an error-free, open-source standard for the global civil engineering community.
\end{quote}

\section{Academic Citation}
To cite HydroAj in academic publications, please use:
\begin{lstlisting}[language=TeX]
@software{bhattarai2026hydroaj,
  author    = {Bhattarai, Ajay},
  title     = {{HydroAj: High-Performance Water Resources Engineering Framework}},
  year      = {2026},
  publisher = {GitHub},
  version   = {0.1.1},
  url       = {https://github.com/ajaybhattarai-123/hydroaj}
}
\end{lstlisting}

\end{document}

